\section{Domain Vocabulary}\label{sec:DomainVocabulary}
If your project was designed using \bdq{Domain-Driven Design}\index{domain-driven design} \autocite{Evans:DomainDrivenDesign}, and you will adjust this Coding Conventions for your project, you should consider to add your project's vocabulary here.

Such a glossary may look like the sample structure below that consist from a table that lists the terms together with a brief description, serving as a kind of index to the terms, and a section that describes the terms in details. If that section gets larger, you may consider to move it to the \nameref{sec:Appendices}.

If you want to add your stuff here, look for the files
\begin{itemize}
\item{\verb#naming/DomainGlossary.tex#}
\item{\verb#naming/DomainGlossary.tbl.tex#}
\end{itemize} 
in the \TeX/\LaTeX\index{TeX}\index{LaTeX} sources for this book on GitHub\index{GitHub} \autocite{TQUADRAT_ORG_DOCUMENT_REPOSITORY}.

%------------------------------------------------------------------------------

\clearpage
\subsection{The Ubiquitous Language Glossary} 
\renewcommand{\cellalign}{tl}
\LTXtable{\linewidth}{naming/DomainGlossary.tbl.tex}

\subsubsection{The Terms}

% The suggested structure for a Domain Glossary

% If you avoid references to other parts of the Coding Conventions document, you can use this
% file also as part of a standalone Domain Glossary, using DomainGlossary.tbl.tex as the
% index to it.

\noindent\rule{\textwidth}{0.4pt}

\paragraph{[\textit{Term Name}]}\label{ddterm:TermName}

\subparagraph{Bounded Context:} [Which context this definition applies to]

\subparagraph{Definition:} 
[Clear, precise definition as agreed upon by domain experts and dev team]

\subparagraph{Synonyms/Aliases:} 
[Other names used for this concept, if any – flag if inconsistent]

\subparagraph{Type:} 
[Entity | Value Object | Aggregate | Aggregate Root | Domain Event | Domain Service | Policy | Specification | Enum/Category]

\subparagraph{Attributes/Properties:} 
[Key attributes, if it's a model element]

\subparagraph{Related Terms:} 
[Links to other glossary entries this term relates to or depends on]

\subparagraph{Example:} 
[Concrete example in plain business language]

\subparagraph{Code Reference:} 
[Class/interface name in the codebase – keeps code and language in sync]

\subparagraph{Notes/Constraints:} 
[Business rules, invariants, edge cases]

\noindent\rule{\textwidth}{0.4pt}

%==============================================================================
% End of file!!
%==============================================================================
